\begin{helper}
Let \(p,m\) be natural numbers and let
\[
  ab:\{0,\ldots,m\}\to \mathbb F_2^p\times\mathbb F_2^p.
\]
Form the length \(m+1\) tuple obtained by keeping the first \(m\) pairs of
\(ab\) and replacing the final pair by \(((ab_m)_1,0)\). Then the rank of the
span of all first coordinates in this modified tuple is equal to the rank of
the span of all first coordinates in \(ab\).
\end{helper}
\begin{proof}
The rank in \(\noderef{F2BadTupleRank}\) depends only on the set of first
coordinates of the tuple entries with index less than \(m+1\). We compare
these two sets directly. For an index \(j<m+1\), either \(j=m\) or \(j<m\).
If \(j=m\), the modified tuple has final first coordinate \((ab_m)_1\), the
same first coordinate as \(ab\). If \(j<m\), the modified tuple agrees with
\(ab\) at \(j\), because its prefix is exactly the restriction of \(ab\) to
the first \(m\) indices. Thus every first coordinate occurring in the modified
tuple occurs in \(ab\), and the same two cases show the reverse inclusion.
The two spans are therefore spans of the same set, so their dimensions, and
hence the two ranks, are equal.
\end{proof}
